%══════════════════════════════════════════════════════════════════════
%  CHAPTER 3 — Proposed Method
%══════════════════════════════════════════════════════════════════════
\chapter{Proposed Method}
\label{ch:proposed}

This chapter presents the proposed \padonap{} (\textit{Proactive Attack Defense over ONAP}) system. It begins with the design rationale and approach direction (Section~\ref{sec:approach}), proceeds through the four-stage pipeline architecture (Section~\ref{sec:arch-overview}), formalises the threat model and safety invariants (Sections~\ref{sec:threat-model}--\ref{sec:safety}), and details the AI components, policy engine, and SLA allocator (Sections~\ref{sec:features}--\ref{sec:payload}).

%%─────────────────────────────────────────────────────────────────────
\section{Approach Direction}
\label{sec:approach}
%%─────────────────────────────────────────────────────────────────────

The central design insight motivating \padonap{} is as follows.
Any reactive DDoS mitigation system has a hard latency floor equal to
the VNF boot time: the orchestrator cannot act faster than its slowest
required function. For a traffic scrubber this floor is approximately
6~s. No amount of orchestration tuning lowers this floor.

The only way to reduce \emph{effective} mitigation latency is to
pre-position the mitigation VNF \emph{before} the attack reaches the
reactive trigger. This requires a forecast. The forecast trades reduced
activation latency (VNF already running when attack peaks) against a
modest false-positive risk (the VNF consumes resources on benign windows
when the forecast was incorrect).

The \padonap{} design addresses this trade-off with three architectural
decisions:

\begin{enumerate}
  \item \textbf{Two-tier proactive action.} Pre-positioning uses
        \tier{2} (rate-limiter, 505~ms boot) rather than \tier{3}
        (scrubber, 6{,}006~ms boot). If the forecast is wrong, the
        cost is 500~ms wasted boot time and moderate resource consumption.
        If correct, the rate-limiter is already in place when the attack
        peaks.

  \item \textbf{LP SLA allocator at every tier transition.} Activating
        any VNF subtracts overhead from the available link capacity.
        An explicit LP solve at each tier change ensures that all tenant
        SLA floors are respected regardless of which tier is active.

  \item \textbf{Hysteresis guard.} A hold window of $k{=}2$ prevents
        rapid oscillation between tiers (VNF churn), which would add
        latency variance and increase resource overhead.
\end{enumerate}

%%─────────────────────────────────────────────────────────────────────
\section{System Architecture Overview}
\label{sec:arch-overview}
%%─────────────────────────────────────────────────────────────────────

\padonap{} is a four-stage pipeline. Every 5-second window, stage~S1
scrapes per-device gNMI counters and NetFlow v5 records;
stage~S2 aggregates them into a 17-dimensional feature vector;
stage~S3 runs the AI inference layer; stage~S4 consumes the AI payload, evaluates
the five-tier policy, allocates tenant bandwidth under an LP SLA model,
and instantiates VNFs through an ONAP SO client.

The DMaaP message bus between S3 and S4 is the key architectural boundary. Because all inter-stage
communication passes through the canonical JSON schema on a named Kafka
topic, any stage can be upgraded or replaced independently.

%%─────────────────────────────────────────────────────────────────────
\section{Threat Model and Safety Invariants}
\label{sec:threat-safety}
%%─────────────────────────────────────────────────────────────────────

We formalise the adversary capability and the system's safety requirements.
The adversary controls a botnet and may emit volumetric floods, protocol-exhaustion attacks, or gradual-ramp attacks. We assume telemetry integrity and control-plane trust.

Three formal safety properties are required:
\begin{itemize}
    \item \textbf{Pre-positioning safety:} For every proactive transition, the bandwidth allocation to every benign tenant must satisfy $a_u(T_j) \geq \text{floor}_u$.
    \item \textbf{URLLC floor invariant:} For every mitigation tier $T$ and in every scenario, $a_{\text{URLLC}}(T) \geq 150$~Mbps.
    \item \textbf{Bounded rollback:} Proactive escalations are held no longer than $k=2$ windows beyond the last supporting signal.
\end{itemize}

%%─────────────────────────────────────────────────────────────────────
\section{AI Components and Feature Engineering}
\label{sec:ai-comp}
%%─────────────────────────────────────────────────────────────────────

The AI layer (Stage S3) embeds two tracks: Track~A, an XGBoost seven-class detector, and Track~B, a Transformer+LSTM four-horizon forecaster.

\paragraph{Feature Engineering.} Every 5-second window is aggregated into a 17-dimensional feature vector $\mathbf{x} \in \mathbb{R}^{17}$ comprising rate features (6), entropy features (5), and protocol-mix features (6). Shannon entropy $H(P) = -\sum p_i \log_2 p_i$ is computed for destination ports, source IPs, and flow segments.

\paragraph{Track A: XGBoost Detector.} The detector uses 400 trees with FGSM adversarial augmentation ($\varepsilon = 0.02$). It identifies attack classes (UDP, SYN, Amplification, HTTP) with high accuracy and low inference latency (2.2~ms).

\paragraph{Track B: Transformer+LSTM Forecaster.} A sliding window of 8 time steps feeds a Transformer encoder and LSTM to produce attack-probability estimates at $t+30, 60, 90, 120$ seconds. This enables the proactive pre-positioning mechanism.

%%─────────────────────────────────────────────────────────────────────
\section{Policy Engine and SLA Allocation}
\label{sec:policy-sla}
%%─────────────────────────────────────────────────────────────────────

The mitigation action space is discretised into five tiers: $T_0$ (Normal), $T_1$ (Alert), $T_2$ (Preempt), $T_3$ (Mitigate), and $T_4$ (Severe).

\paragraph{LP SLA Allocator.} At tier $T$ with VNF overhead $o(T)$, the allocator solves:
\begin{align}
  \max_{\mathbf{a}} &\quad \sum w_u \, a_u \nonumber \\
  \text{s.t.} &\quad \sum a_u \;\leq\; C - o(T), \nonumber \\
  &\quad f_u \;\leq\; a_u \;\leq\; d_u, \nonumber
\end{align}
guaranteeing per-tenant floors $f_u$ (e.g., 150~Mbps for URLLC).

%%─────────────────────────────────────────────────────────────────────
\section{Chapter Summary}
\label{sec:proposed-summary}
%%─────────────────────────────────────────────────────────────────────

This chapter presented the full \padonap{} design, combining a four-stage pipeline architecture with a dual-track AI engine and an LP-backed SLA allocator. The proactive approach direction addresses the structural latency of reactive VNF instantiation by leveraging multi-horizon forecasting for pre-positioning. The system ensures safety through formal invariants and a graduated five-tier policy ladder. The next chapter evaluates this proposed method across eight orchestration scenarios.
